\section{Analysis of Model-$A$ ($\gamma_1=\gamma_2=0$, $\theta_1=\theta_2=0$)}\label{sec:model_a}

Model-$A$ is the baseline of this collection of models: all jockeying and abandonment parameters are set to zero, leaving an ordinary two-class $M/M/1$ queueing model with non-preemptive priority. The analysis is carried out in both state spaces $S$ and $\widetilde{S}$: on the former, two complete proofs are presented; while on the latter, an approximation and its corresponding evaluation are provided.

The balance equations are given —Equations \eqref{eq:gen:balance_interior}, \eqref{eq:gen:balance_x_boundary}, \eqref{eq:gen:balance_y_boundary}, \eqref{eq:gen:balance_zero} and \eqref{eq:gen:balance_idle} for state space $S$, and Equations \eqref{eq:gen:tilde_balance_interior}, \eqref{eq:gen:tilde_balance_diagonal}, \eqref{eq:gen:tilde_balance_y_boundary}, \eqref{eq:gen:tilde_balance_zero} and \eqref{eq:gen:tilde_balance_idle} for state space $\widetilde{S}$ by plugging $\gamma_1=\gamma_2=0$ and $\theta_1=\theta_2=0$. The derivation also holds accordingly, so the starting points are the fundamental equations presented in \ref{lem:gen:fundamental_equation} for $S$ and \ref{lem:gen:PPGF_dynamics} for $\widetilde{S}$. The corresponding subsections will introduce the equation for their dynamics, with all jockeying and abandonments parameters set to zero.

The stability conditions, since we have no abandonments, will be the same as for an ordinary sincle-class $M/M/1$ queueing system, \ref{eq:A:stability}: the system will be positive recurrent if and only if
\begin{equation*}
    \rho = \frac{\lambda_1+\lambda~_2}{\mu} <1.
\end{equation*}

\textbf{Stability.} With no abandonment, the system is positive recurrent if and only if the total traffic intensity satisfies
\begin{equation}
    \rho = \frac{\lambda_1+\lambda_2}{\mu} < 1.
    \label{eq:A:stability}
\end{equation}
The lack of abandonments gives us
\begin{equation}
    \pi_0= 1-\rho \quad \text{and} \quad \pi(0,0)=(1-\rho) \rho.
\end{equation}

The figure below corresponds to the queueing diagram of Model-$A$.
\begin{figure}[H]
    \centering
    \resizebox{0.62\textwidth}{!}{%
        % Model A: γ₁=γ₂=θ₁=θ₂=0
% Two queues, one server. Arrivals and service only — no jockeying, no abandonment.
\begin{tikzpicture}[>=Latex, thick]

    % ── Queue 1 (top, Class-1, priority) ──────────────────────────────────────
    \draw (0, 1.9) -- (2, 1.9) -- (2, 1.1) -- (0, 1.1);
    \foreach \x in {0.2, 0.4, ..., 1.8} { \draw[thin] (\x, 1.8) -- (\x, 1.2); }
    \draw[->] (-1.5, 1.5) -- (-0.1, 1.5) node[midway, above] {$\lambda_1$};

    % ── Queue 2 (bottom, Class-2) ──────────────────────────────────────────────
    \draw (0, -1.1) -- (2, -1.1) -- (2, -1.9) -- (0, -1.9);
    \foreach \x in {0.2, 0.4, ..., 1.8} { \draw[thin] (\x, -1.2) -- (\x, -1.8); }
    \draw[->] (-1.5, -1.5) -- (-0.1, -1.5) node[midway, above] {$\lambda_2$};

    % ── Server ────────────────────────────────────────────────────────────────
    \node[draw, circle, minimum size=1.2cm] (server) at (5, 0) {$\mu$};
    \draw[->] (2.2,  1.5) -- (server);
    \draw[->] (2.2, -1.5) -- (server);
    \draw[->] (server) -- (7, 0);

\end{tikzpicture}
%
    }
    \caption{Queue layout for Model-$A$. Only Poisson arrivals and exponential service
    are present; all jockeying and abandonment rates are zero.}
    \label{fig:model-a-queue}
\end{figure}

\subsection{Analysis model A: state space $S$}
We start by stating Equation \eqref{eq:gen:fundamental_equation} with no jockeying and abandonment parameters, where the marginal $P_y(y)=P(0,y)$ is still to be determined.
\begin{align}
    \begin{split}
        \left[ (\lambda_1+\lambda_2+ \mu)xy - \mu y - \lambda_1 x^2 y - \lambda_2 x y^2 \right] P(x,y) 
         &= \mu(x-y) P_y(y) 
         - \mu x(1-y) \pi(0,0).
    \end{split} \label{eq:A:fundamental}
\end{align}

The following theorem summarizes the main result of this subsection, and two standalone proofs will be provided: one analytical and one using probabilistic arguments.
\begin{theorem}\label{thm:A:PGF}
    Consider a queueing system with $2$ priority classes with Poisson arrival rates $\lambda_1$ and $\lambda_2$, where class-$1$ customers have non-preemptive priority over those of class-$2$, and exponentially distributed service times with rate $\mu$ independently of the class.
 The PGF $P(x,y)$ of such a system is given by
    \begin{align}
        P(x,y) &= \left[ (1+\rho_1 +\rho_2)xy - y -\rho_1x^2y - \rho_2xy^2 \right]^{-1} \cdot y(1-y) \left[ \frac{x - x^*(y)}{x^*(y)-y} \right] \cdot \rho (1-\rho) \\[10pt]
        &\text{where} \nonumber \\[5pt]
        x^*(y) &= \frac{\left[\mu + \lambda_1 + \lambda_2 (1-y)\right] - \sqrt{\left[\mu + \lambda_1 + \lambda_2 (1-y)\right]^2 - 4 \lambda_1 \mu}}{2 \lambda_1}
    \end{align}
\end{theorem}

\begin{corollary}\label{cor:A:pi0}
    Under the stability condition $\rho<1$, the empty-state probabilities for Model-$A$ are
    \begin{equation}
        \pi_0 = 1-\rho, \qquad \pi(0,0) = \rho(1-\rho). \label{eq:A:pi0_pi00}
    \end{equation}
    In particular, $P(1,1) = 1-\pi_0 = \rho$ and the diagonal generating function satisfies $P(z,z)=\pi(0,0)/(1-\rho z)$.
\end{corollary}

\begin{proof}[Proof of Corollary~\ref{cor:A:pi0}]
    \textbf{PGF of the total number of customers}  Setting $x=y=z$ in the fundamental equation (Equation \eqref{eq:A:fundamental}):
    \begin{align*}
        \begin{split}
            \left[ (\lambda_1+\lambda_2+ \mu)z^2 - \mu z - \lambda_1 z^3 - \lambda_2 z^3 \right] P(z,z) 
            &= \mu(z-z) P_y(z)  - \mu z(1-z) \pi(0,0)
        \end{split} \\
        \begin{split}
            \left[ (\lambda_1+\lambda_2+ \mu)z - \mu  - (\lambda_1 + \lambda_2) z^2 \right] P(z,z) 
            &= - \mu (1-z) \pi(0,0)
        \end{split} \\
        \begin{split}
            -(1-z)\left[ \mu - (\lambda_1 + \lambda_2)z \right] P(z,z) 
            &= - \mu (1-z) \pi(0,0)
        \end{split} \\
        \begin{split}
            \left[ \mu - (\lambda_1 + \lambda_2)z \right] P(z,z) 
            &= \mu \pi(0,0)
        \end{split} \\
        \begin{split}
            P(z,z) 
            &= \frac{\mu \pi(0,0)}{\mu - \lambda z}
        \end{split} \\
        \begin{split}
            P(z,z) &= \frac{\pi(0,0)}{1-\rho} 
        \end{split}
    \end{align*}
    Using the normalisation condition for the PGF, $P(1,1) = \rho$, we can solve for $\pi(0,0)$ and obtain 
    \begin{equation*}
        \pi(0,0) = (1 - \rho) \rho,
    \end{equation*}
    and with $\lambda\pi_0=\mu \pi(0,0)$, we recover
    \begin{equation*}
        \pi_0=1-\rho.
    \end{equation*}
\end{proof}

\subsubsection{Analytical approach of model A in state space $S$}

\begin{proof}

    We start by statimg again the fundamental equation for Model-$A$ (Equation \eqref{eq:A:fundamental}):
    \begin{equation*}
        \left[ (\lambda_1+\lambda_2+ \mu)xy - \mu y - \lambda_1 x^2y - \lambda_2 xy^2 \right] P(x,y) 
        = \mu(x-y) P_y(y)  - \mu x(1-y) \pi(0,0).
    \end{equation*}
    We apply the \emph{Kernel Method}: we fix $y\in(0,1]$ and find $x^*(y)$ such that the LHS the expression above is zero. Given that $P(x,y)$ is a PGF, it is bounded. At any such root, the LHS is zero; therefore the RHS must also be zero: i.e. $(x^*(y)-y)\,P_y(y) \;=\; x^*(y)(1-y)\,\pi(0,0)$.
    To find the roots of the coefficient of $P(x,y)$, we set
    \begin{align*}
        \begin{split}
             0 &=  (\lambda_1 + \lambda_2 + \mu)xy - \mu y -\lambda_1x^2y - \lambda_2xy^2
        \end{split} \\
        \begin{split}
             0 &=  (\lambda_1 + \lambda_2 + \mu)x - \mu  -\lambda_1x^2 - \lambda_2xy
        \end{split} \\
        \begin{split}
             0 &=  - \lambda_1 x^2 + \left[ \mu + \lambda_1 + \lambda_2 (1-y) \right]x - \mu 
        \end{split} \\
        \begin{split}
          0 &=  \lambda_1 x^2 - \left[ \mu + \lambda_1 + \lambda_2 (1-y) \right] x + \mu 
        \end{split} \\
        \begin{split}
            x^*(y) &= \frac{
                \left[\mu + \lambda_1 + \lambda_2 (1-y) \right] \pm \sqrt{
                    \left[\mu + \lambda_1 + \lambda_2 (1-y) \right]^2 - 4  \mu \lambda_1
                }
            }{
                2\lambda_1
            }
        \end{split}
    \end{align*}
    We label the root with the negative sign $x^*(y)$ and the root with the positive sign $x^+(y)$. At $y=1$, we have
    \begin{equation*}
        x^*(1) = 1 \quad \text{and} \quad x^+(1) > 1
    \end{equation*}
    We need our root to be inside the unit disk in the interval $[0,1)$. Taking the derivative of $x^*(y)$, we see that ,
    \begin{equation}
        \frac{d}{dy}x^*(y) = \frac{\lambda_2}{2\lambda_1}\!\left[-1+\frac{\mu+\lambda_1+\lambda_2(1-y)}{\sqrt{(\mu+\lambda_1+\lambda_2(1-y))^2-4\lambda_1\mu}}\right]>0,
        \label{eq:xstar_prime}
    \end{equation}
    since $(\mu+\lambda_1+\lambda_2(1-y))^2>(\mu+\lambda_1+\lambda_2(1-y))^2-4\lambda_1\mu$ whenever $\lambda_1\mu>0$. Therefore, $x^*(y)<1$ for $x\in[0,1)$. By means of Vieta's formulas for the kernel polynomial $\lambda_1 x^2-[\mu+\lambda_1+\lambda_2(1-y)]x+\mu=0$, we have that the product of the two roots is $x^*(y) x^+(y)= \mu/\lambda_1>1$. Therefore, ($x^+(y)=\mu/(\lambda_1 x^*(y))>1$ throughout $[0,1)$, so the only root that is strictly inside the unit disk for the interval $[0,1)$ is $x^*(y)$. We now state
    \begin{equation}
        x^*(y) = \frac{
            (\mu +\lambda_1 + \lambda_2 (1-y)) 
            - \sqrt{(\mu +\lambda_1 + \lambda_2 (1-y))^2 - 4 \lambda_1 \mu}
        }{
            2\lambda_1
        }.
        \label{eq:A:xstar}
    \end{equation}
    Since the next steps will involve L'Hôpital's rule in order to compute a limit at $y\to 1$, we recall $$\frac{d}{dy}x^*(y)$$ from Equation \eqref{eq:xstar_prime} and we evaluate at $y=1$ to see how it behaves:
    \begin{align*}
        \frac{d}{dy}x^*(1)
        &= \frac{\lambda_2}{2 \lambda_1}
        \left[ 
            -1 + \frac{
                \mu + \lambda_1
            }{
                \sqrt{(\mu+\lambda_1)^2- 4\lambda_1\mu}
            }
        \right]\\
        &= \frac{\lambda_2}{2\lambda_1}
        \left[ 
            -1 + \frac{
                \mu + \lambda_1
            }{
                \sqrt{(\mu-\lambda_1)^2}
            }
        \right]\\
        &= \frac{\lambda_2}{2\lambda_1}
        \left[
            \frac{\mu + \lambda_1 - (\mu - \lambda_1)}{\mu - \lambda_1}
        \right] \\
        &= \frac{\lambda_2}{\mu - \lambda_1} \\
        &= \frac{\rho_2}{1-\rho_1}
    \end{align*}
    For future reference,
    \begin{equation}
        \frac{d}{dy}x^*(1) = \frac{\rho_2}{1-\rho_1}.
        \label{eq:A:xstar_deriv}
    \end{equation}

    Next, back to the fundamental equation for Model-$A$ (Equation~\eqref{eq:A:fundamental}), its RHS at $x=x^*(y)$ also equals zero: we have $\mu \big[ (x-y) P_y(y) - x(1-y) \pi(0,0) \big] = 0$, so
    \begin{equation}
        P_y(y) = \frac{x^*(y)(1-y)}{x^*(y) - y} \pi(0,0). \label{eq:A:Py}
    \end{equation}
    
    From Corollary~\ref{cor:A:pi0}, we know that $\pi(0,0)=(1-\rho)\rho$. Substituting this and Equation~\eqref{eq:A:Py} simultaneously into our fundamental equation (Equation~\eqref{eq:A:fundamental}) yields:
    \begin{align*}
        \begin{split}
            \left[ (\lambda_1+\lambda_2+ \mu)xy - \mu y - \lambda_1 x^2y - \lambda_2 xy^2 \right] P(x,y) 
            &= \mu(x-y) \frac{x^*(y)(1-y)}{x^*(y) - y} \rho(1-\rho)  - \mu x(1-y) \rho(1-\rho)
        \end{split} \\
        \begin{split}
            \left[ (\lambda_1+\lambda_2+\mu)xy -\mu y-\lambda_1 x^2y - \lambda_2 xy^2 \right] P(x,y)
            &= \mu (1-y) \left[ \frac{x^*(y)(x-y)}{x^*(y)-y} - x \right] \rho (1-\rho)
        \end{split} \\
        \begin{split}
            \left[ (\lambda_1+\lambda_2+\mu)xy -\mu y-\lambda_1 x^2y - \lambda_2 xy^2 \right] P(x,y)
            &= \mu  (1-y) \left[ \frac{y(x - x^*(y))}{x^*(y) - y} \right] \rho (1-\rho)
        \end{split} \\
        \begin{split}
            \left[ (1 + \rho_1 + \rho_2)xy - y-\rho_1 x^2y - \rho_2 xy^2 \right] P(x,y)
            &= y (1-y) \left[ \frac{x - x^*(y)}{x^*(y) - y} \right] \rho (1-\rho)
        \end{split}
    \end{align*}
    
    Multiplying both sides by $\left[(1 + \rho_1 + \rho_2)xy - y-\rho_1 x^2y - \rho_2 xy^2\right]^{-1}$ to isolate $P(x,y)$ concludes the proof.                              
\end{proof}




\subsubsection{Probabilistic approach of model A in state space $S$}
\begin{proof}
    Our goal here is to derive the partial PGF $P_y(y)$ based on probabilistic arguments. As done in the analytical approach, we start from the fundamental global balance equation for Model-$A$, i.e.\ Equation \eqref{eq:A:fundamental}:
    \begin{equation}
        \left[ (\lambda_1 + \lambda_2 + \mu) xy - \mu y - \lambda_1 x^2 y - \lambda_2 x y^2 \right] P(x,y) =
        \mu \big[ (x-y) P_y(y) - x(1-y) \pi(0,0) \big].
        \label{eq:A:prob_starting}
    \end{equation}
    
    Here, $P_y(y) = P(0,y) = \sum_{n_2=0}^{\infty} \pi(0,n_2)y^{n_2}$ represents the partial PGF restricted to states where the server is busy and the high-priority queue is empty ($N_1 = 0$). From the definition of conditional expectation, we can express $P_y(y)$ in the following way:
    \begin{equation}
        P_y(y) = \mathbb{P}(\text{busy}, N_1=0) \cdot \mathbb{E}[y^{N_2} \mid \text{busy}, N_1=0].
        \label{eq:A:Py_prob_decomp}
    \end{equation}
    
    In order to address the conditional expectation of class-$2$ customers, we adopt their perspective. Since class-$1$ has non-preemptive priority, class-$2$ customers are blocked whenever $N_1>0$; their effective service time is therefore the duration of a complete busy period $B$ of the $M/M/1(\lambda_1,\mu)$ sub-queue, so from class-$2$'s viewpoint the system is an $M/G/1(\lambda_2, B)$ queue. The LST and mean of $B$ are given in~\S\ref{ssec:mm1} (Equations~\eqref{eq:mm1:bp_lst}--\eqref{eq:mm1:bp_mean}) with $\lambda\leftarrow\lambda_1$:
    \begin{align}
        \widetilde{B}(s) &= \frac{(\mu+\lambda_1+s) - \sqrt{(\mu + \lambda_1 + s)^2-4\mu\lambda_1}}{2\lambda_1}, \label{eq:A:B_lst}\\
        \mathbb{E}[B] &= \frac{1}{\mu(1-\rho_1)}. \label{eq:A:B_mean}
    \end{align}

    In this $M/G/1$ framework, we apply the Pollaczek--Khinchine (PK) formula (\S\ref{ssec:pk}, Equation~\eqref{eq:pk:pgf}) to find $L_2(y)$, the PGF of the number of customers in the equivalent $M/G/1$ system:
    \begin{equation}
        L_2(y) = \frac{(1 - \lambda_2 \mathbb{E}[B]) (1-y) \widetilde{B}(\lambda_2(1-y))}{\widetilde{B}(\lambda_2(1-y)) - y}.
        \label{eq:A:pk_formula}
    \end{equation}

    Evaluating $\widetilde{B}(s)$ at $s = \lambda_2(1-y)$ recovers precisely the kernel root $x^*(y)$ from Equation~\eqref{eq:A:xstar}:
    \begin{equation}
        \widetilde{B}(\lambda_2(1-y)) = \frac{\mu+\lambda_1+\lambda_2(1-y) - \sqrt{(\mu + \lambda_1 + \lambda_2(1-y))^2-4\mu\lambda_1}}{2\lambda_1} = x^*(y).
    \end{equation}

    Substituting $\widetilde{B}(\lambda_2(1-y))=x^*(y)$ and $\lambda_2 \mathbb{E}[B] = \rho_2/(1-\rho_1)$ into Equation~\eqref{eq:A:pk_formula} gives:
    \begin{equation}
        L_2(y) = \frac{1-\rho}{1-\rho_1} \cdot \frac{x^*(y)(1-y)}{x^*(y) - y}.
        \label{eq:A:L2_pgf}
    \end{equation}

    By the PASTA property, class-$2$ Poisson arrivals see the time-average joint distribution. Since the class-$2$ subsystem is isomorphic to an $M/G/1(\lambda_2,B)$ queue — consecutive class-$2$ service initiations are separated by i.i.d.\ intervals $B$, and class-$2$ customers in this sub-queue correspond exactly to $N_2$ waiting class-$2$ customers in the original model when $N_1=0$ — the time-average PGF of $N_2$ conditional on $(N_1=0, \text{busy})$ equals the $M/G/1$ steady-state PGF:
    \begin{equation}
        \mathbb{E}[y^{N_2} \mid \text{busy}, N_1=0] = L_2(y) = \frac{1-\rho}{1-\rho_1} \cdot \frac{x^*(y)(1-y)}{x^*(y) - y}.
        \label{eq:A:L2_conditional}
    \end{equation}

    For the prefactor, we determine $\mathbb{P}(\text{busy},N_1=0)$ by a level-crossing argument on the class-$1$ waiting queue. Under non-preemptive priority, whenever $N_1\geq1$ the server always selects a waiting class-$1$ customer upon any service completion. Therefore the flux crossing level $N_1=n_1$ upward (a class-$1$ arrival when $N_1=n_1$, rate $\lambda_1$) equals the flux downward (a service completion removes one waiting class-$1$, rate $\mu$):
    \[
        \lambda_1\,\mathbb{P}(N_1=n_1-1,\,\text{busy}) = \mu\,\mathbb{P}(N_1=n_1,\,\text{busy}), \quad n_1\geq1.
    \]
    This recursion gives $\mathbb{P}(N_1=n_1,\,\text{busy})=\rho_1^{n_1}\,\mathbb{P}(N_1=0,\,\text{busy})$.  Summing over all $n_1\geq0$ and using $\mathbb{P}(\text{busy})=\rho$:
    \[
        \mathbb{P}(N_1=0,\,\text{busy})\cdot\frac{1}{1-\rho_1}=\rho \implies \mathbb{P}(\text{busy},N_1=0)=\rho(1-\rho_1).
    \]
    In particular, $\mathbb{P}(N_1=0\mid\text{busy})=1-\rho_1$.

    Returning to Equation~\eqref{eq:A:Py_prob_decomp}:
    \begin{align*}
         &P_y(y) \\
         & \quad=\quad \mathbb{P}(\text{busy}, N_1=0) \cdot \mathbb{E}[y^{N_2} \mid \text{busy}, N_1=0] \\
         & \quad=\quad \rho(1-\rho_1) \cdot \left( \frac{1-\rho}{1-\rho_1} \cdot \frac{x^*(y)(1-y)}{x^*(y) - y} \right)\\
         &\quad=\quad \rho (1-\rho) \cdot \frac{x^*(y)(1-y)}{x^*(y) - y}. 
    \end{align*}
    
    We now substitute $P_y(y)$ and $\pi(0,0)=\rho(1-\rho)$ back into the right-hand side of our starting global balance Equation \eqref{eq:A:prob_starting}:
    \begin{align*}
        \left[ (\lambda_1 + \lambda_2 + \mu) xy - \mu y - \lambda_1 x^2 y - \lambda_2 x y^2 \right] P(x,y) 
        &= \mu \left[ (x-y) \left( \rho(1-\rho) \frac{x^*(y)(1-y)}{x^*(y) - y} \right) - x(1-y) \rho(1-\rho) \right].
    \end{align*}
    
    Factoring out the common terms $\mu \rho (1-\rho) (1-y)$ gives:
    \begin{align*}
        &\left[ (\lambda_1 + \lambda_2 + \mu) xy - \mu y - \lambda_1 x^2 y - \lambda_2 x y^2 \right] P(x,y) \\
        &\quad=\quad \mu (1-y) \left[ \frac{(x-y)x^*(y)}{x^*(y) - y} - x \right] \rho(1-\rho)\\
        &\quad=\quad  \mu (1-y) \left[ \frac{x \cdot x^*(y) - y \cdot x^*(y) - x(x^*(y) - y)}{x^*(y) - y} \right] \rho(1-\rho)\\
        &\quad=\quad  \mu (1-y) \left[ \frac{xy - y \cdot x^*(y)}{x^*(y) - y} \right]\rho(1-\rho) \\
        &\quad=\quad  \mu y(1-y) \left[ \frac{x - x^*(y)}{x^*(y) - y} \right]\rho(1-\rho).
    \end{align*}
    
    Dividing both sides by $\mu$ yields the following:
    \begin{equation*}
        \left[ (1 + \rho_1 + \rho_2) xy - y - \rho_1 x^2 y - \rho_2 x y^2 \right] P(x,y) = y(1-y) \left[ \frac{x - x^*(y)}{x^*(y) - y} \right]\rho(1-\rho).
    \end{equation*}
    
    And, finally, multiplying both sides by $\left[ (1+\rho_1 +\rho_2)xy - y -\rho_1x^2y - \rho_2xy^2 \right]^{-1}$ completes the proof:
    \begin{equation}
        P(x,y) = \left[ (1+\rho_1 +\rho_2)xy - y -\rho_1x^2y - \rho_2xy^2 \right]^{-1}  y(1-y) \left[ \frac{x - x^*(y)}{x^*(y)-y} \right] \rho (1-\rho).
    \end{equation}
\end{proof}




\subsubsection{PPGF approach of model A in state space $S$}

In this subsection, we approach the balance equations by keeping the class-$1$ index $N_1$ as a discrete parameter and forming the PGF only with respect to $N_2$, defining $\widetilde{P}(n_1,y)=\sum_{n_2\geq0}\pi(n_1,n_2)\,y^{n_2}$. Multiplying the interior and boundary balance equations by $y^{n_2}$ and summing yields a second-order linear recurrence in $n_1$ whose characteristic roots determine the geometric structure of the solution in $n_1$.

From the interior balance equation, then, we obtain
\begin{align*}
    \begin{split}
       & (\lambda_1 + \lambda_2 + \mu) \sum_{n_2=1}^{\infty} \pi(n_1,n_2) y^{n_2} \\
        &\quad = \quad\mu \sum_{n_2=1}^{\infty} \pi(n_1+1,n_2) y^{n_2}
        + \lambda_1 \sum_{n_2=1}^{\infty} \pi(n_1-1,n_2) y^{n_2} 
        + \lambda_2 \sum_{n_2=1}^{\infty} \pi(n_1,n_2-1) y^{n_2}
    \end{split} \\
    \begin{split}
        & (\lambda_1 + \lambda_2 + \mu) \left[ \sum_{n_2=0}^{\infty} \pi(n_1,n_2) y^{n_2} - \pi(n_1,0) \right] \\
        & \quad = \quad\mu \left[ \sum_{n_2=0}^{\infty} \pi(n_1+1,n_2) y^{n_2} - \pi(n_1+1,0) \right]\\
        &\quad\quad+ \lambda_1 \left[ \sum_{n_2=0}^{\infty} \pi(n_1-1,n_2) y^{n_2} - \pi(n_1-1,0) \right]\\
        &\quad\quad+ \lambda_2 y \sum_{n_2=0}^{\infty} \pi(n_1,n_2) y^{n_2} 
    \end{split} \\
    \begin{split}
        &(\lambda_1 + \lambda_2 + \mu) \left[ \widetilde{P}(n_1,y) - \pi(n_1,0) \right] \\
        & \quad= \quad \mu \left[ \widetilde{P}(n_1+1, y) - \pi(n_1+1,0)\right]
        + \lambda_1 \left[ \widetilde{P}(n_1-1,y) - \pi(n_1-1,0) \right]
        + \lambda_2 y \widetilde{P}(n_1,y)
    \end{split} 
\end{align*}

Note that we can use the $x$-boundary equation (Equation \eqref{eq:gen:balance_x_boundary}) to simplify the expression. Doing so and then grouping terms yields
\begin{equation}
    \left[ \mu + \lambda_1 + \lambda_2 (1-y) \right] \widetilde{P}(n_1,y)
    = \mu \widetilde{P}(n_1+1,y)
    + \lambda_1 \widetilde{P}(n_1-1,y).
    \label{eq:A:ppgf_recurrence}
\end{equation}

Since Equation~\eqref{eq:A:ppgf_recurrence} is a second-order linear recurrence in $n_1$ with constant (in $n_1$) coefficients, its general solution is a linear combination of geometric sequences. Because $\widetilde{P}(n_1,y)\to 0$ as $n_1\to\infty$ (the queue cannot grow without bound under stability), only one geometric mode survives, giving
\begin{equation}
    \widetilde{P}(n_1, y) = \widetilde{P}(0,y)\,[\widetilde{x}^*(y)]^{n_1},
\end{equation}
where $\widetilde{x}^*(y)$ is the root with $|\widetilde{x}^*(y)|\leq 1$ of the characteristic equation
\begin{equation}
    \mu [\widetilde{x}^*(y)]^2 - [\mu + \lambda_1 + \lambda_2(1-y)] \widetilde{x}^*(y) + \lambda_1 = 0.
\end{equation}

Solving the characteristic equation gives two roots:
\begin{equation*}
    \widetilde{x}^*(y) = \frac{(\mu + \lambda_1 + \lambda_2(1-y)) \pm \sqrt{(\mu + \lambda_1 + \lambda_2(1-y))^2 - 4\lambda_1\mu}}{2\mu}.
\end{equation*}

We select the root with the negative sign. At $y=1$ the characteristic polynomial factors as $(\mu \widetilde{x} - \lambda_1)(\widetilde{x} - 1) = 0$, giving roots $\rho_1<1$ and $1$. The root equal to $1$ would make the series $\sum_{n_1=0}^{\infty} \widetilde{P}(0,y)[\widetilde{x}^*(y)]^{n_1}$ diverge, contradicting $P(x,1)<\infty$ for $|x|<1$. Hence only the root with the negative sign, which satisfies $|\widetilde{x}^*(y)|<1$ for $|y|\leq 1$, is admissible:
\begin{equation}
    \widetilde{x}^*(y) = \frac{(\mu + \lambda_1 + \lambda_2(1-y)) - \sqrt{(\mu + \lambda_1 + \lambda_2(1-y))^2 - 4\lambda_1\mu}}{2\mu}.
\end{equation}

In particular, at $y=1$ the root recovers the class-$1$ traffic intensity:
\begin{equation}
    \widetilde{x}^*(1) = \frac{\lambda_1}{\mu} = \rho_1,
\end{equation}
consistent with the geometric marginal $P(N_1=n_1,\text{busy})=(1-\rho_1)\rho_1^{n_1}\rho$ of the priority class. To complete this approach: substitute $\widetilde{P}(n_1,y)=\widetilde{P}(0,y)[\widetilde{x}^*(y)]^{n_1}$ into the $y$-boundary equation~\eqref{eq:gen:balance_y_boundary} to determine $\widetilde{P}(0,y)$, then reconstruct
\[
    P(x,y)=\sum_{n_1=0}^\infty\widetilde{P}(0,y)[\widetilde{x}^*(y)]^{n_1}x^{n_1}=\frac{\widetilde{P}(0,y)}{1-\widetilde{x}^*(y)x};
\]
the resulting formula should match Theorem~\ref{thm:A:PGF}. These final steps are not carried out here.

\subsection{Analysis of Model~$A$: state space $\widetilde{S}$}

Setting $\gamma_1=\gamma_2=\theta_1=\theta_2=0$ in the $\widetilde{S}$ balance equations~\eqref{eq:gen:tilde_balance_interior}--\eqref{eq:gen:tilde_balance_idle} and forming the PPGF $\widetilde{P}(y,n)=\sum_{n_2=0}^n \widetilde{\pi}(n_2,n)\,y^{n_2}$ yields the inhomogeneous recurrence:
\begin{align}
    \begin{split}
       &\left[ \lambda_1 + \lambda_2 + \mu \right] \widetilde{P}(y,n) \\
       & \quad = \quad(\lambda_1 + \lambda_2 y) \widetilde{P}(y,n-1) + \mu \widetilde{P}(y,n+1) + \mu y^n (1-y) \tpi(n+1,n+1)
    \end{split} \label{eq:A:ppgf_Stilde_dynamics}
\end{align}

\begin{theorem}[Approximate PPGF on $\widetilde{S}$; heuristic]\label{thm:A:PPGF_Stilde}
    \emph{(The following result rests on an unproved decay assumption; see the proof.)} Under the approximation that neglects the boundary-diagonal forcing term $\mu y^n(1-y)\,\tpi(n+1,n+1)$ in~\eqref{eq:A:ppgf_Stilde_dynamics}, the homogeneous recurrence admits the geometric solution
    \begin{equation}
        \widetilde{P}(y,n) \;\approx\; \rho(1-\rho)\,[\widetilde{y}^*(y)]^n, \label{eq:A:Ptilde_approx}
    \end{equation}
    where
    \begin{equation}
        \widetilde{y}^*(y) = \frac{(\lambda+\mu) - \sqrt{(\lambda+\mu)^2 - 4\mu(\lambda_1+\lambda_2 y)}}{2\mu}, \qquad \lambda = \lambda_1+\lambda_2. \label{eq:A:ytildestar}
    \end{equation}
    At $y=1$ this recovers $\widetilde{P}(1,n)=(1-\rho)\rho^{n+1}$, consistent with the $M/M/1(\lambda,\mu)$ steady-state marginal for $n$ customers waiting.
\end{theorem}

\begin{proof}[Proof of Theorem~\ref{thm:A:PPGF_Stilde}]
For state space $\widetilde{S}$, the dynamics for the PPGF are described by the following inhomogeneous difference equation:
\begin{equation}
    (\lambda_1 + \lambda_2 + \mu) \widetilde{P}(y,n) = (\lambda_1 + \lambda_2 y) \widetilde{P}(y,n-1) + \mu \widetilde{P}(y,n+1) + \mu y^n (1-y) \tpi(n+1,n+1).
\end{equation}
The inhomogeneous term $\mu y^n(1-y)\tpi(n+1,n+1)$ prevents the direct application of Cohen's geometric-sequence ansatz $\tP(y,n)=\tP(y,0)\cdot [f(y)]^n$. The diagonal probabilities $\tpi(n,n)=\pi(0,n)$ correspond to states where all $n$ waiting customers are class-$2$; from the exact formula in Theorem~\ref{thm:A:PGF} one can verify that $\pi(0,n)$ decays geometrically in $n$, but this is not derived here.  \emph{Assuming} this geometric decay, the forcing term is exponentially small and we obtain an approximation by neglecting it. The resulting homogeneous recurrence is
\begin{equation}
    (\lambda_1 + \lambda_2 + \mu) \widetilde{P}(y,n) = (\lambda_1 + \lambda_2 y) \widetilde{P}(y,n-1) + \mu \widetilde{P}(y,n+1).
\end{equation}
By equating it to $0$, we can obtain the characteristic polynomial of this difference equation:
\begin{equation}
    \mu [\widetilde{y}(y)]^2 - (\lambda_1 + \lambda_2 + \mu) \widetilde{y}(y) + (\lambda_1 + \lambda_2 y) = 0,
\end{equation}
where solving for $\widetilde{y}(y)$ yields
\begin{equation}
    \widetilde{y}^*(y) = \frac{(\lambda_1 + \lambda_2 + \mu) \pm \sqrt{(\lambda_1 + \lambda_2 + \mu)^2 - 4\mu(\lambda_1 + \lambda_2 y)}}{2\mu}.
\end{equation}
Plugging in $y=1$ already lets us exclude the positive root —$\widetilde{y}_{+}(y)=1$ and $\widetilde{y}_{-}(y)=\rho$— and leaves us with
\begin{equation}
    \widetilde{y}^*(y) = \frac{(\lambda_1 + \lambda_2 + \mu) - \sqrt{(\lambda_1 + \lambda_2 + \mu)^2 - 4\mu(\lambda_1 + \lambda_2 y)}}{2\mu}.
\end{equation}
Applying Cohen's trick yields
\begin{equation}
    \tP(y,n)=\tP(y,0) \cdot  [\widetilde{y}^*(y)]^n.
\end{equation}
Using $\tP(y,0)=\sum_{n_2=0}^0\tpi(n_2,0)y^{n_2}=\tpi(0,0)=(1-\rho)\rho$, we obtain
\begin{equation}
    \tP(y,n)=(1-\rho)\rho \cdot  [\widetilde{y}^*(y)]^n.
\end{equation}
We can evaluate this expression at $y=1$ and check that it will behave like a regular $M/M/1$ queueing system with arrival rate $\lambda=\lambda_1+\lambda_2$ and service rate $\mu$. Note that, since $n$ counts the total number of customers in the queues and not in the system, we have an exponent $n+1$ in the $\rho$ term, with that $+1$ accounting for the customer in service:
\begin{equation}
    \tP(1,n)= (1-\rho)\rho \cdot  [\widetilde{y}^*(1)]^n = (1-\rho)\rho \cdot\rho^n= (1-\rho) \rho^{n+1}. \qedhere
\end{equation}
\end{proof}

\subsection{Mean queue lengths and mean waiting times}

\subsubsection*{Class-1 marginal from the theorem formula}

We evaluate $P(x,y)$ of the theorem at $y=1$ for fixed $x\neq1$.
The denominator satisfies $D(x,y)\to D(x,1) = -\rho_1(x-1)(x-1/\rho_1)$,
while the factor $x-x^*(y)\to x-x^*(1)=x-1$.  For the remaining $0/0$,
expand $x^*(y)$ near $y=1$:
\begin{align}
    \begin{split}
        x^*(y) \;=\; 1 - \frac{\rho_2}{1-\rho_1}(1-y) + O\!\bigl((1-y)^2\bigr),
    \end{split}
\end{align}
so $x^*(y)-y = \frac{1-\rho}{1-\rho_1}(1-y)+O((1-y)^2)$ and the factor
$(1-y)$ in numerator and denominator cancel:
\begin{align}
    \begin{split}
        P(x,1)
        &\;=\; \lim_{y\to1} \frac{y\,\rho(1-\rho)\,(x-x^*(y))}{D(x,y)\,(x^*(y)-y)/(1-y)} \\
        &\;=\; \frac{\rho(1-\rho)\cdot(x-1)}{-\rho_1(x-1)(x-1/\rho_1)\cdot\frac{1-\rho}{1-\rho_1}} \\
        &\;=\; \frac{\rho(1-\rho_1)}{-\rho_1(x-1/\rho_1)}
         \;=\; \frac{\rho(1-\rho_1)}{1-\rho_1 x}.
    \end{split}
    \label{eq:A:Px1}
\end{align}
Differentiating at $x=1$:
\begin{align}
    \begin{split}
        \mathbb{E}[N_1] \;=\; \left.\frac{d}{dx}P(x,1)\right|_{x=1}
        \;=\; \frac{\rho(1-\rho_1)\,\rho_1}{(1-\rho_1)^2}
        \;=\; \frac{\rho\rho_1}{1-\rho_1}.
    \end{split}
    \label{eq:A:EN1}
\end{align}

\subsubsection*{Total queue length from the diagonal of the theorem formula}

Observe that $x^*(z)=z$ iff $z$ is a root of $\lambda(z-1)(z-1/\rho)=0$, so
$x^*(z)\neq z$ for all $z\in(0,1)$ with $\rho<1$.  Therefore the ratio
$(z-x^*(z))/(x^*(z)-z)=-1$ throughout $(0,1)$, and the theorem formula
reduces to:
\begin{align}
    \begin{split}
        P(z,z)
        \;=\; \frac{z(1-z)\,\rho(1-\rho)}{D(z,z)}\cdot(-1).
    \end{split}
\end{align}
The denominator at $x=y=z$ is
$D(z,z)=(1+\rho)z^2-z-\rho z^3=-\rho z(z-1)(z-1/\rho)$,
so:
\begin{align}
    \begin{split}
        P(z,z)
        &\;=\; \frac{-z(1-z)\,\rho(1-\rho)}{-\rho z(z-1)(z-1/\rho)} \\
        &\;=\; \frac{(1-z)(1-\rho)}{(z-1)(z-1/\rho)}
         \;=\; \frac{-(1-\rho)}{z-1/\rho}
         \;=\; \frac{\rho(1-\rho)}{1-\rho z}.
    \end{split}
    \label{eq:A:Pzz}
\end{align}
Since $\frac{d}{dz}P(z,z)\big|_{z=1}
= \frac{\partial P}{\partial x}(1,1)+\frac{\partial P}{\partial y}(1,1)
=\mathbb{E}[N_1]+\mathbb{E}[N_2]$:
\begin{align}
    \begin{split}
        \mathbb{E}[N_1]+\mathbb{E}[N_2]
        \;=\; \left.\frac{d}{dz}\frac{\rho(1-\rho)}{1-\rho z}\right|_{z=1}
        \;=\; \frac{\rho^2}{1-\rho}.
    \end{split}
    \label{eq:A:EN}
\end{align}
This is the $M/M/1(\lambda,\mu)$ queue-length result for the combined system.

\subsubsection*{Class-2 mean by subtraction}

\begin{align}
    \begin{split}
        \mathbb{E}[N_2]
        &\;=\; \frac{\rho^2}{1-\rho} - \frac{\rho\rho_1}{1-\rho_1}
         \;=\; \frac{\rho^2(1-\rho_1)-\rho\rho_1(1-\rho)}{(1-\rho)(1-\rho_1)} \\
        &\;=\; \frac{\rho(\rho-\rho_1)}{(1-\rho)(1-\rho_1)}
         \;=\; \frac{\rho\rho_2}{(1-\rho_1)(1-\rho)}.
    \end{split}
    \label{eq:A:EN2}
\end{align}

\subsubsection*{Mean waiting times via Little's Law}

\textbf{Little's Law} applied to the \emph{waiting sub-system} of each class: since $N_i$ counts only customers \emph{waiting} in queue (excluding the customer in service), Little's Law gives $\mathbb{E}[N_i]=\lambda_i\,\mathbb{E}[W_i]$ where $\mathbb{E}[W_i]$ is the mean waiting time in queue (time from arrival until service start, not including the service duration itself):
\begin{align}
    \begin{split}
        \mathbb{E}[W_1] \;=\; \frac{\mathbb{E}[N_1]}{\lambda_1}
        \;=\; \frac{\rho\rho_1}{\lambda_1(1-\rho_1)}
        \;=\; \frac{\rho}{\mu(1-\rho_1)},
    \end{split}
    \label{eq:A:EW1} \\
    \begin{split}
        \mathbb{E}[W_2] \;=\; \frac{\mathbb{E}[N_2]}{\lambda_2}
        \;=\; \frac{\rho\rho_2}{\lambda_2(1-\rho_1)(1-\rho)}
        \;=\; \frac{\rho}{\mu(1-\rho_1)(1-\rho)}.
    \end{split}
    \label{eq:A:EW2}
\end{align}
Class-1 gains from priority: $\mathbb{E}[W_1]$ does not depend on $\rho_2$,
while $\mathbb{E}[W_2]$ carries the extra factor $1/(1-\rho_1)$.








